\chapter{The PyObject Model and Reference Counting (Python 1.x)}

Chapter 1 established that a Python program is a graph of heap objects and that names are
bindings into it. This chapter dissects the node of that graph. \textbf{Every Python value
is a heap-allocated C struct that begins with the same header: a reference count and a type
pointer.} From those two fields flow polymorphism (the type pointer), automatic memory
management (the count), and the entire object protocol. We also confront the two facts that
most surprise engineers measuring Python today: an integer costs 28 bytes, and
\texttt{None}'s reference count is over four billion and never changes. Both follow directly
from the model.

\section*{Section Index}
\begin{itemize}
  \item \textbf{2.1} The universal object header: \texttt{PyObject} and \texttt{PyVarObject}
  \item \textbf{2.2} The type object and the slot system
  \item \textbf{2.3} Reference counting: the lifecycle, and what \texttt{Py\_INCREF}/\texttt{Py\_DECREF} really do
  \item \textbf{2.4} Immortal objects (PEP 683) and what \texttt{getrefcount} now reports
  \item \textbf{2.5} What reference counting cannot do: cycles
  \item \textbf{2.6} Caching, constant deduplication, and interning---three different mechanisms
  \item \textbf{2.7} Performance and memory: the cost of universal boxing
  \item \textbf{2.8} Anti-patterns and the senior-engineer contrast
  \item \textbf{2.9} Summary and cross-references
\end{itemize}

\section{The universal object header: \texttt{PyObject} and \texttt{PyVarObject}}

\textbf{Why this exists.} For the interpreter to treat every value uniformly---store it in a
list, pass it to a function, garbage-collect it, print it---every value must expose the same
minimal interface. CPython prefixes \emph{every} object with a common header. That header is
the contract that makes ``everything is an object'' mechanically true.

\begin{lstlisting}[language=C, caption={Illustrative; Include/object.h. PyObject is the base of every Python value.}]
typedef struct _object {
    _PyObject_HEAD_EXTRA          /* debug-only doubly-linked list (Py_TRACE_REFS builds) */
    Py_ssize_t    ob_refcnt;      /* reference count */
    PyTypeObject *ob_type;        /* pointer to the type that defines behavior */
} PyObject;

typedef struct {
    PyObject    ob_base;
    Py_ssize_t  ob_size;          /* number of items in the variable part */
} PyVarObject;
\end{lstlisting}

Two fields, and everything else is built on them: \texttt{ob\_refcnt} drives memory
management (§2.3), and \texttt{ob\_type} drives behavior (§2.2). A fixed-size object embeds
\texttt{PyObject} directly; a variable-size object (\texttt{list}, \texttt{str},
arbitrary-precision \texttt{int}) embeds \texttt{PyVarObject} and stores its payload after
the header.

\textbf{The senior-engineer contrast.} In C++ an object's layout is fixed by its static type
at compile time, and virtual dispatch goes through a per-class \textbf{vtable} pointer.
CPython inverts this: a value's size depends on its \emph{runtime} type while every
\emph{reference} is one pointer wide (8 bytes on 64-bit); \emph{all} operations dispatch
through the per-object \texttt{ob\_type} pointer; and a type is itself a first-class runtime
object (\texttt{type(x)} is a \texttt{PyObject}).

\begin{lstlisting}[caption={Memory layout of a boxed value}]
name x --> +-------------------------------------+  heap
           | (debug link fields, debug builds)   |
           | ob_refcnt : 1                       |
           | ob_type   : --> PyTypeObject 'int'  |
           | <payload: the integer's digits>     |
           +-------------------------------------+
\end{lstlisting}

Because all references are pointer-width, a \texttt{list} of a million integers is a million
pointers (8 MB) pointing at a million separate 28-byte integer objects scattered on the
heap---not a packed array of machine words. That fact is why \texttt{array}, \texttt{numpy},
and the buffer protocol (Vol IX) exist; §2.7 quantifies the cost.

\section{The type object and the slot system}

\textbf{Why this exists.} \texttt{ob\_type} points to a \texttt{PyTypeObject}---itself an
object---holding the function pointers (``slots'') the interpreter calls to make an object
\emph{do} things. When you write \texttt{a + b}, the interpreter reads \texttt{a}'s type's
\texttt{tp\_as\_number->nb\_add} slot, a direct C function pointer; it does not search for a
method named \texttt{\_\_add\_\_} by string at that moment. The dunder protocol you write is
the \emph{source}; the slots are the \emph{compiled mechanism}.

\begin{lstlisting}[language=C, caption={Illustrative; Include/cpython/object.h. PyTypeObject is itself a PyVarObject.}]
struct _typeobject {
    PyObject_VAR_HEAD
    const char *tp_name;                /* "module.Class" for repr/errors */
    Py_ssize_t  tp_basicsize, tp_itemsize;
    destructor  tp_dealloc;             /* called when refcount hits 0 (sec 2.3) */
    reprfunc    tp_repr;                /* __repr__ */
    hashfunc    tp_hash;               /* __hash__ */
    newfunc     tp_new;                /* __new__  : allocate */
    initproc    tp_init;               /* __init__ : initialize */
    allocfunc   tp_alloc;              /* low-level allocation */
    PyNumberMethods   *tp_as_number;   /* nb_add, nb_multiply, ... */
    PySequenceMethods *tp_as_sequence; /* sq_item, sq_length, ... */
    PyMappingMethods  *tp_as_mapping;  /* mp_subscript, ... */
};
\end{lstlisting}

\textbf{Dunder-to-slot mapping.} When a class is created, CPython wires each dunder into its
slot: \texttt{\_\_init\_\_}$\rightarrow$\texttt{tp\_init},
\texttt{\_\_new\_\_}$\rightarrow$\texttt{tp\_new},
\texttt{\_\_repr\_\_}$\rightarrow$\texttt{tp\_repr},
\texttt{\_\_hash\_\_}$\rightarrow$\texttt{tp\_hash},
\texttt{\_\_del\_\_}$\rightarrow$\texttt{tp\_finalize} (not \texttt{tp\_dealloc}),
\texttt{\_\_add\_\_}$\rightarrow$\texttt{tp\_as\_number->nb\_add},
\texttt{\_\_getitem\_\_}$\rightarrow$\texttt{tp\_as\_mapping->mp\_subscript} or
\texttt{tp\_as\_sequence->sq\_item}.

\begin{lstlisting}[language=Python, caption={Defining \_\_hash\_\_ fills the tp\_hash slot; the builtin hash() calls it.}]
class CustomObject:
    def __init__(self, val):
        self.val = val
    def __hash__(self):          # -> tp_hash
        return hash(self.val)

print("hash matches delegate:", hash(CustomObject("test")) == hash("test"))
\end{lstlisting}

Verified output (CPython 3.13.5):

\begin{lstlisting}[caption={Output}]
hash matches delegate: True
\end{lstlisting}

Grouping numeric/sequence/mapping methods into sub-tables is a space optimization: a
non-numeric type leaves \texttt{tp\_as\_number} \texttt{NULL} rather than carrying dozens of
unused pointers. The full slot system and descriptor protocol are the subject of Vol VIII;
here we need only that \textbf{behavior lives on the type, reached through
\texttt{ob\_type}.}

\section{Reference counting: the lifecycle, and what \texttt{Py\_INCREF}/\texttt{Py\_DECREF} really do}

\textbf{Why this exists.} CPython's primary memory manager is \textbf{reference counting}:
each object counts how many references point at it; when the count reaches zero the object is
freed \emph{immediately and deterministically}. This contrasts with the tracing collectors of
Java, Go, and the CLR, and it has concrete consequences: deterministic finalization and no
stop-the-world pause for acyclic garbage, but per-operation counting overhead and a hard
problem with cycles (§2.5).

Every reference gained calls \texttt{Py\_INCREF}; every reference lost calls
\texttt{Py\_DECREF}; the NULL-safe variants \texttt{Py\_XINCREF}/\texttt{Py\_XDECREF} guard
null pointers. When a decref brings the count to zero, the \textbf{deallocation pipeline}
runs: (1) the interpreter calls \texttt{obj->ob\_type->tp\_dealloc(obj)}; (2) the
deallocator \textbf{decrefs every object this object referenced}---freeing a list decrefs
each element, freeing a dict decrefs every key and value, a recursive depth-first cascade;
(3) the raw memory returns to the allocator (PyMalloc---Vol VIII).

\begin{lstlisting}[language=Python, caption={Reference counting is observable; subtract 1 for getrefcount's own argument ref.}]
import sys

my_list = [100, 200]
print("getrefcount-1 at start:", sys.getrefcount(my_list) - 1)
alias = my_list
print("after one alias:", sys.getrefcount(my_list) - 1)
del alias
print("after del alias:", sys.getrefcount(my_list) - 1)
\end{lstlisting}

Verified output (CPython 3.13.5):

\begin{lstlisting}[caption={Output}]
getrefcount-1 at start: 1
after one alias: 2
after del alias: 1
\end{lstlisting}

\texttt{del alias} does \emph{not} delete the list; it removes one \emph{binding}. The object
dies only when the last reference goes away---the operational meaning of ``names are
bindings.''

\textbf{Implementation detail, not language semantics.} Reference counting is a CPython
mechanism. PyPy, GraalPy, and Jython use tracing collectors and do \emph{not} refcount; code
relying on a \texttt{\_\_del\_\_} firing at a precise moment relies on CPython, not on Python.
Use \texttt{with}/context managers for deterministic release (Chapter 5).

\section{Immortal objects (PEP 683) and what \texttt{getrefcount} now reports}

\textbf{Why this exists.} A handful of objects---\texttt{None}, \texttt{True},
\texttt{False}, the small integers, the empty \texttt{str}/\texttt{tuple}, interned
identifiers, type objects---are referenced from everywhere and live for the whole process.
Counting their references is pure overhead, and in a \textbf{free-threaded} build (Vol VII)
every such write is a contended atomic on a shared cache line. \textbf{PEP 683 (Python 3.12)}
made these objects \textbf{immortal}: their reference count is pinned to a sentinel, and
\texttt{Py\_INCREF}/\texttt{Py\_DECREF} skip them.

\begin{lstlisting}[language=Python, caption={Immortal objects (3.12+) report a pinned, enormous refcount that never changes.}]
import sys
print("refcount(None):", sys.getrefcount(None))
print("as hex:", hex(sys.getrefcount(None)))
\end{lstlisting}

Verified output (CPython 3.13.5):

\begin{lstlisting}[caption={Output}]
refcount(None): 4294967295
as hex: 0xffffffff
\end{lstlisting}

That \texttt{0xffffffff} is the immortality sentinel on this build, not a real count. Never
branch on the \emph{magnitude} of a refcount; the modern \texttt{Py\_INCREF} macro first
checks for the immortal sentinel; and the overhead this removes is what makes the
free-threaded build's shared-object access viable (Vol VII, Ch 20).

\section{What reference counting cannot do: cycles}

Reference counting has one fatal blind spot: a \textbf{reference cycle} keeps its own counts
above zero even when nothing outside the cycle can reach it.

\begin{lstlisting}[language=Python, caption={A self-referential structure that pure refcounting can never free.}]
a = {}
b = {}
a["b"] = b
b["a"] = a          # a -> b -> a : each holds the other; neither count reaches 0
del a, b            # bindings gone, but the two dicts still reference each other
\end{lstlisting}

After \texttt{del a, b}, the two dictionaries are unreachable yet each still has refcount 1
(from the other). Pure reference counting would leak them forever. CPython's answer is a
separate, optional \textbf{cyclic garbage collector} (the \texttt{gc} module) that
periodically finds and reclaims such islands---the subject of Chapter 3, where it enters the
story with Python 2.0. Hold the division of labor: \textbf{refcounting frees acyclic garbage
instantly; the cyclic collector cleans up the rest.}

\section{Caching, constant deduplication, and interning---three different mechanisms}

These three are routinely conflated, producing the most common piece of Python
misinformation: ``\texttt{257 is 257} is \texttt{False}.'' Run as a script, \textbf{it is
\texttt{True}.} Three distinct mechanisms are at work.

\textbf{The small-integer cache (runtime singleton pool).} CPython preallocates the integers
from $-5$ to $256$ as shared singletons. Any expression producing a value in that range
returns the same object, even across separate code objects.

\textbf{Constant deduplication (compile-time, per code object).} When the compiler builds a
code object it stores each distinct literal once in \texttt{co\_consts} and deduplicates
equal constants. Two \texttt{257} literals in the same compilation unit resolve to the same
object---even though 257 is outside the small-int cache. This is not interning.

\textbf{String interning (runtime global table).} Identifier-like string literals are added
to a process-global \textbf{intern table}; equal identifier-like literals share one object.
Strings built at runtime are not auto-interned; \texttt{sys.intern()} forces it.

\begin{lstlisting}[language=Python, caption={Three mechanisms, demonstrated and distinguished.}]
import sys

a = 256; b = 256
print("256 cached (small-int pool):", a is b)

c = 257; d = 257
print("257 in one code object (constant dedup):", c is d)

g1 = {}; exec("x = 257", g1)
g2 = {}; exec("x = 257", g2)
print("257 across two compilation units:", g1["x"] is g2["x"])

s1 = "godhood"; s2 = "godhood"
print("identifier-like literal interned:", s1 is s2)

r2 = "".join(["god", "hood"])      # built at runtime
print("runtime-built string auto-interned:", r2 is s1)
print("sys.intern() forces sharing:", sys.intern(r2) is s1)

t1 = (); t2 = tuple()
print("empty tuple singleton:", t1 is t2)
\end{lstlisting}

Verified output (CPython 3.13.5):

\begin{lstlisting}[caption={Output}]
256 cached (small-int pool): True
257 in one code object (constant dedup): True
257 across two compilation units: False
identifier-like literal interned: True
runtime-built string auto-interned: False
sys.intern() forces sharing: True
empty tuple singleton: True
\end{lstlisting}

Corrected model: \texttt{257 is 257} is \texttt{True} in a script (constant dedup) and
\texttt{False} only when the two \texttt{257}s come from different compilation units---which
is exactly what the interactive REPL does, compiling each line separately. That is the source
of the folklore.

\textbf{Anti-pattern.} Never use \texttt{is} to compare values. \texttt{is} is identity; its
results depend on caching, dedup, and interning, none of which are language guarantees for
general values. Use \texttt{==}. Reserve \texttt{is} for genuine singletons (\texttt{is
None}, sentinels). Modern CPython emits a \texttt{SyntaxWarning} for \texttt{is} applied
directly to a literal.

\section{Performance and memory: the cost of universal boxing}

Universal boxing---every value a separate heap object behind a pointer---makes the data model
uniform and is also Python's dominant performance tax.

\begin{lstlisting}[language=Python, caption={The memory cost of boxed objects (bytes), CPython 3.13.5, 64-bit.}]
import sys
for label, obj in [("int 0", 0), ("int 2**30", 2**30), ("int 2**100", 2**100),
                   ("list []", []), ("list [1,2,3]", [1, 2, 3]),
                   ("tuple ()", ()), ("str ''", ""), ("str 'a'", "a")]:
    print(f"{label:14s} {sys.getsizeof(obj):3d} bytes")
\end{lstlisting}

Verified output (CPython 3.13.5):

\begin{lstlisting}[caption={Output}]
int 0           28 bytes
int 2**30       32 bytes
int 2**100      40 bytes
list []         56 bytes
list [1,2,3]    88 bytes
tuple ()        40 bytes
str ''          41 bytes
str 'a'         42 bytes
\end{lstlisting}

\begin{itemize}
  \item \textbf{An integer is 28 bytes}, not 8---a \texttt{PyVarObject} with a header plus
    arbitrary-precision digits, growing with magnitude. A list of a million ints is
    $\sim$8 MB of pointers \emph{plus} $\sim$28 MB of integer objects, versus 8 MB for a C
    array. This factor, plus lost cache locality from chasing scattered pointers, is the
    motivation for \texttt{array.array}, \texttt{numpy}, and the buffer protocol (Vol IX).
  \item \textbf{Refcount churn is real CPU cost.} Every bind, argument pass, and append
    touches a refcount. Under the GIL these are plain increments; in the free-threaded build
    they become atomics---hence immortalization (§2.4) and biased reference counting
    (Vol VII).
  \item \textbf{Identity caching saves allocations} for hot tiny objects, which is why the
    model is viable despite per-object overhead.
\end{itemize}

\section{Anti-patterns and the senior-engineer contrast}

\begin{itemize}
  \item \textbf{\texttt{is} for value comparison} (§2.6): a latent bug that ``works'' until
    the value leaves the cached range or the REPL. Use \texttt{==}.
  \item \textbf{Relying on \texttt{\_\_del\_\_} for cleanup}: refcount-driven finalization is
    a CPython detail, not guaranteed promptly, and never runs for cycle members until the
    cyclic collector acts. For OS resources use a context manager---Python's closest analog
    to C++ RAII (Chapter 5).
  \item \textbf{Assuming value semantics on assignment} (the C++ reflex): \texttt{b = a}
    aliases; it does not copy. Use \texttt{copy.copy}/\texttt{copy.deepcopy} for
    independence.
  \item \textbf{Treating \texttt{sys.getsizeof} as deep size}: it returns the object's own
    size, not the size of referenced objects.
\end{itemize}

\section{Summary and cross-references}

\begin{itemize}
  \item Every value is a heap object beginning with a shared header: \textbf{\texttt{ob\_refcnt}}
    (memory) and \textbf{\texttt{ob\_type}} (behavior); variable-length objects add
    \texttt{ob\_size}.
  \item Behavior lives on the \textbf{type object} as C function-pointer \textbf{slots}; the
    dunder methods you write fill those slots.
  \item \textbf{Reference counting} frees acyclic garbage deterministically, at the cost of
    per-operation counting and an inability to reclaim \textbf{cycles} ($\rightarrow$
    Chapter 3).
  \item \textbf{PEP 683 immortal objects} (3.12+) pin process-global refcounts;
    \texttt{getrefcount(None)} reports \texttt{0xffffffff}.
  \item Identity surprises come from \textbf{three separate mechanisms}---the small-int
    cache, compile-time constant deduplication, and string interning. \texttt{257 is 257} is
    \texttt{True} in a script. Compare values with \texttt{==}, never \texttt{is}.
  \item Universal boxing makes an \texttt{int} cost 28 bytes and destroys array locality---the
    root motivation for the numeric/zero-copy machinery in Vol IX.
\end{itemize}

\textbf{Cross-references.} The cyclic garbage collector $\rightarrow$ Chapter 3. Context
managers and deterministic cleanup $\rightarrow$ Chapter 5. The full type-slot and descriptor
machinery $\rightarrow$ Vol VIII. The free-threaded refcounting model $\rightarrow$ Vol VII,
Ch 20. \texttt{array}, \texttt{memoryview}, and zero-copy numerics $\rightarrow$ Vol IX.
